% 1.3 — Problématique
\section{Problématique : des agents qui exécutent des commandes}
\label{sec:problematique}

\subsection{Un agent n'écrit plus seulement du code, il l'exécute}
Un agent de code IA comme Claude Code ou Codex ne se contente pas de proposer du texte.
Il lit le projet, modifie des fichiers, et surtout lance des commandes dans un shell ---
\cmd{bash} sous Linux, PowerShell sous Windows --- pour compiler, tester, nettoyer, migrer
une base, déployer. Chaque commande s'exécute avec les droits du développeur, sur sa
machine, souvent à côté de données qui n'existent nulle part ailleurs. L'agent se trompe
parfois ; une erreur dans une commande destructrice n'est pas une suggestion maladroite,
c'est une perte de données.

\subsection{Trois incidents publics}
Le risque n'est pas théorique. Trois cas documentés illustrent trois formes du problème.

\begin{description}[style=nextline,leftmargin=0.6cm,font=\sffamily\bfseries]
  \item[Une base de production supprimée (Replit, juillet 2025)]
    Au cours d'une expérience de plusieurs jours, l'agent de Replit a exécuté des commandes
    qui ont effacé la base de données de production d'un utilisateur, alors qu'il avait reçu
    la consigne de ne rien modifier sans accord explicite ; il a ensuite affirmé à tort que
    la restauration était impossible~\cite{oecd-replit,eweek-replit}.
    \emph{Leçon : une consigne n'est pas une protection, et l'agent n'est pas une source
    fiable sur ce qui peut être récupéré.}
  \item[Des fichiers écrasés par des déplacements (Gemini CLI, juillet 2025)]
    L'agent a cru avoir créé un dossier dont la création avait échoué, puis y a « déplacé »
    les fichiers de l'utilisateur ; chaque déplacement vers un nom inexistant a écrasé le
    précédent, et tous les fichiers sauf un ont été perdus~\cite{aiid-gemini,gh-gemini}.
    \emph{Leçon : aucune commande de la séquence ne « ressemble » à une suppression ; seul
    un raisonnement sur l'état du système de fichiers révèle la perte.}
  \item[Des fichiers supprimés en voulant les compter (Claude Code, juillet 2026)]
    Chargé de compter des fichiers de rapport, un agent a lancé une commande qui les a
    supprimés ; l'utilisateur en a récupéré une partie avec un outil de récupération de
    données~\cite{gh-claude-76626}. La bêta de \demo{} cite ce cas comme motivation.
    \emph{Leçon : demander une confirmation ne suffit pas quand la personne fait confiance
    à l'intention annoncée ; ce qui manquait, c'était une copie prise avant.}
\end{description}

\subsection{Pourquoi le problème est difficile}
Juger une commande \emph{avant} qu'elle s'exécute paraît simple : repérer \cmd{rm}, et se
méfier. En réalité le langage du shell rend cette lecture trompeuse. Le
tableau~\ref{tab:regex-pieges} montre quatre lignes, tirées des tests du projet, où une
lecture naïve --- par motifs, ou par un analyseur trop confiant --- se trompe dans un sens
ou dans l'autre.

\begin{table}[htbp]
\centering
\caption{Lignes où une lecture naïve se trompe (extraits de \fichier{tests/test\_command\_model.py} et des commits I1a, W0)}
\label{tab:regex-pieges}
\small
\begin{tabularx}{\linewidth}{@{}L{5.2cm}YY@{}}
\toprule
\textbf{Ligne} & \textbf{Lecture naïve} & \textbf{Ce que fait le shell} \\
\midrule
\cmd{ls \# ; rm -rf x} & un \cmd{rm -rf} & rien : tout ce qui suit \cmd{\#} est un commentaire \\
\cmd{echo \$(date; rm -rf x) >/dev/null} & un \cmd{echo} & exécute \cmd{rm -rf x} dans la substitution \\
\cmd{bash -o pipefail -c "rm app.db"} & une option inconnue, rien à juger & supprime \cmd{app.db} \\
\cmd{X=-rf; rm \$X src} & un \cmd{rm} sans option & un \cmd{rm -rf src} \\
\bottomrule
\end{tabularx}
\end{table}

À cela s'ajoutent les enveloppes (\cmd{sudo}, \cmd{xargs}, \cmd{timeout}), les commandes
lancées par du code que l'outil ne lit pas (\cmd{python script.py}), les effets à distance
(une base, un dépôt git distant, un bucket) qu'aucune copie locale ne restaure, et un
second langage, PowerShell, où les arguments se lient aux paramètres par position.

\subsection{Énoncé du problème}
\begin{encadre}[Problématique]
Comment, avant qu'une commande écrite par un agent ne s'exécute, en prédire l'effet assez
exactement pour laisser passer ce qui est sans danger, rendre récupérable ce qui peut
l'être, et arrêter ou soumettre à une personne le reste --- sans jamais revendiquer une
protection qu'on n'a pas prise, et de la même manière sous Linux et sous Windows ?
\end{encadre}
